% Generated draft: requires mathematical and source review.
\documentclass[11pt]{article}
\usepackage[margin=25mm]{geometry}
\usepackage{fontspec}
\setmainfont{Latin Modern Roman}
\usepackage{amsmath,amssymb,mathtools,unicode-math}
\setmathfont{Latin Modern Math}
\usepackage{enumitem,needspace,xurl,hyperref}
\setlength{\parindent}{0pt}
\setlength{\parskip}{5pt}
\setlength{\emergencystretch}{3em}
\newcommand{\N}{\mathbb{N}}
\newcommand{\Z}{\mathbb{Z}}
\newcommand{\Q}{\mathbb{Q}}
\newcommand{\R}{\mathbb{R}}
\newcommand{\C}{\mathbb{C}}
\newcommand{\F}{\mathbb{F}}
\newcommand{\A}{\mathbb{A}}
\newcommand{\PP}{\mathbb{P}}
\newcommand{\E}{\mathbb{E}}
\newcommand{\eps}{\varepsilon}
\newcommand{\dd}{\,\mathrm{d}}
\begin{document}
\section*{UnsolvedMath ID: 20002672}
\textbf{Source problem code:} catalogue entry
\subsection{Definitions and mathematical statement}
This damaged extraction joins a free-Gibbs compactness question and a separate free-entropy-dimension question. Matrix size and tuple length are denoted differently below to avoid the source's mixed indices.

\textbf{1. Compactness of the limiting free heat semigroup.} Consider a normalizable matrix model with density proportional to \(\exp[-N\operatorname{Tr}P(A_1,A_1^*,\ldots,A_m,A_m^*)]\), and suppose a limiting tracial law exists. The source recalls the conjugate-variable relation \(J_j=D_jP\), with \(D_j\) the cyclic derivative, and asks whether
\[
 e^{-t\mathcal A},\qquad \mathcal A=\sum_{j=1}^m\partial_j^*\partial_j,
\]
is compact on the limiting \(L^2\) space, for positive time. Here \(\partial_j\) are Voiculescu free difference quotients with respect to the limiting generators; the operator must be defined from the closure of the quadratic form \(q(\xi)=\sum_j\|\partial_j\xi\|_2^2\) with the coarse product-trace target. Closability, domains and the self-adjoint realization are necessary qualifications. The source suggests beginning with \(P=\sum_iA_i^2+\sum_\ell t_\ell q_\ell(A)\), citing convergence work of Guionnet and Maurel-Segala, but gives no smallness or convergence hypotheses. The aid concerns the self-adjoint finite-tuple case, not every complex matrix model in the initial display.

For a nonnegative self-adjoint \(\mathcal A\), compactness of \(e^{-t\mathcal A}\) for any, equivalently every, \(t>0\) is equivalent to compact resolvent. The aid reports this property for the finite variance-one semicircular tuple \(S=(S_1,\ldots,S_m)\): its free-gradient Laplacian is the number operator \(\mathcal N\), with eigenvalue \(d\) on degree \(d\) and multiplicity \(m^d\), \(d\ge0\). Its heat trace is finite exactly when \(me^{-t}<1\), and then equals \((1-me^{-t})^{-1}\). Compactness at every positive time does not imply trace class at every positive time.

The aid's transfer criterion assumes a trace-preserving analytic free transport, a Jacobian boundedly invertible on the coarse \(L^2\) bimodule, and compatible cores for the closed gradient forms. Equivalently for the comparison step, its induced \(L^2\) unitary \(U\) must give constants \(0<a\le b<\infty\) such that
\[
 a q_0(\xi)\le q(U\xi)\le b q_0(\xi)
\]
on a common form core, hence on the corresponding closed form domains. The reported consequences are compact resolvent and eigenvalue comparisons
\[
 a\lambda_k(\mathcal N)\le\lambda_k(\mathcal A)
                    \le b\lambda_k(\mathcal N),
\]
with multiplicities, and eventual heat trace class, for example when \(at>\log m\). The aid applies this to sufficiently small analytic self-adjoint free-Gibbs perturbations for which those transport hypotheses hold; no quantitative smallness theorem is supplied in the record.

For a correlated quadratic semicircular law with positive covariance matrix \(C\), the reported exact spectrum consists of all finite sums of precision eigenvalues \(\nu_1,\ldots,\nu_m\) of \(C^{-1}\):
\[
 0,\quad\nu_{i_1}+\cdots+\nu_{i_d},\qquad d\ge1,
\]
counted by words. In this normalization the heat trace is \((1-\sum_i e^{-t\nu_i})^{-1}\) when \(\sum_i e^{-t\nu_i}<1\). The variance-one convention corresponds to potential \(\tfrac12\sum_iA_i^2\); the source's \(\sum_iA_i^2\) has a different precision normalization. Arbitrary free-Gibbs compactness is not settled by these conditional comparisons.

\textbf{2. Reconstructed entropy-dimension question.} Let \(X=(x_1,\ldots,x_m)\) be a bounded self-adjoint tuple and let \(S\) be a variance-one semicircular tuple free from \(X\). Put \(X_t=X+\sqrt t\,S\). The extraction has a damaged second delta symbol, inconsistent indices, and printed limsup operations. The supplied verification explicitly replaces those displays by the following reconstructed standard formulas, rather than treating the OCR as exact:
\[
 \delta^*(X)=m-\liminf_{t\downarrow0}
          \frac{\chi^*(X_t)}{\log\sqrt t},\qquad
 \delta^\star(X)=m-\liminf_{t\downarrow0}t\Phi^*(X_t).
\]
Here \(\chi^*\) is nonmicrostates free entropy and \(\Phi^*\) is the total free Fisher information \(\sum_j\|J_j\|_2^2\), not an additional repeated sum of that total. The question is whether \(\delta^*(X)=\delta^\star(X)\) in general, in these explicitly reconstructed conventions.

Using the de Bruijn relation \(\frac{\dd}{\dd t}\chi^*(X_t)=\tfrac12\Phi^*(X_t)\), define \(u=-\log t\) and \(g(u)=e^{-u}\Phi^*(X_{e^{-u}})\). When the usual regularization and finite-reference-entropy hypotheses justify integration, the aid obtains
\[
 \delta^*(X)=m-\liminf_{T\to\infty}\frac1T\int_0^Tg(u)\,\dd u,
 \qquad
 \delta^\star(X)=m-\liminf_{u\to\infty}g(u).
\]
Thus equality holds whenever \(t\Phi^*(X_t)\) has a limit as \(t\downarrow0\). The remaining issue is equality of the lower logarithmic Ces\`aro average and the pointwise lower limit; the aid does not remove this Tauberian gap in general. It is unrelated to proving compactness for all limiting matrix models.
\subsection{Short English statement}
The merged record asks about compact free-Gibbs heat semigroups and equality of two nonmicrostates entropy dimensions. The aid transfers compact resolvent from a semicircular number operator under explicit analytic-transport and form-core hypotheses. For the separately reconstructed dimension definitions, de Bruijn's identity reduces equality to a logarithmic averaging issue, settled when the Fisher-information scaling has a limit.
\subsection{Sources}
\begin{itemize}
\item Local UnsolvedMath registry, permanent ID 20002672, catalogue entry; canonical \texttt{aim-probability-notes.json}, notes[113].
\item AIM \emph{Free Analysis}, \url{https://aimath.org/WWN/freeanalysis/freeanalysis.pdf}. The extraction crosses the heading ``0.4 Focus Group on Free Entropy (day 3)'' and merges two questions. Source status is reconstructed unverified; the supplied verification explicitly repairs the delta symbols and liminf formulas.
\item Guionnet and Maurel-Segala are named in the source for convergence of perturbative matrix models; Voiculescu is named for free difference quotients and the entropy-dimension framework. No titles, years or theorem references are supplied.
\item Machine-generated research attempt 1, partial result, attempt status open, reported rigor proved and low novelty confidence; registry status partially solved. No authoritative separate citation for the transport refinement or the Tauberian reduction is supplied; the aid is not peer reviewed.
\end{itemize}
\end{document}
